\documentclass{article}
\usepackage[fontset=ubuntu]{ctex}
\usepackage{cube}
\usepackage{enumitem}
\usepackage[most]{tcolorbox}
\usepackage[margin=.5cm,a4paper,landscape]{geometry}
\usepackage{array,tabularx,colortbl}

% region 个性化设计
% region 自定义颜色
\definecolor{title}{RGB}{34, 98, 163} % {222, 244, 255}
\definecolor{fll}{RGB}{203, 223, 244} % {222, 244, 255}
\definecolor{oll}{RGB}{242, 218, 193}
\definecolor{pll}{RGB}{244, 210, 227}
\definecolor{myblue}{RGB}{212, 227, 255} % {203, 223, 244}
\definecolor{steps}{RGB}{243, 251, 240} % {194, 234, 154}
\definecolor{notation}{RGB}{235, 243, 251}
% endregion

% region 自定义命令
\ProvideDocumentCommand{\red}{m}{\textcolor{red}{#1}}
\ProvideDocumentCommand{\blue}{m}{\textcolor{blue}{#1}}
\ProvideDocumentCommand{\orange}{m}{\textcolor{orange}{#1}}
\ProvideDocumentCommand{\green}{m}{\textcolor{nicegreen}{#1}}
% endregion

% region 设置表格线条粗细和内边距
\setlength\arrayrulewidth{0.3pt}  % 竖隔线细一点
\setlength\tabcolsep{2pt}         % 表格内边距为1em
\pagestyle{empty} % 去掉页眉页脚页码
% endregion

% region 中文字体的设置
% 设置 CJK 主字体（对应 \rmfamily）为思源宋体 Light
\setCJKmainfont{Noto Serif CJK SC}[
    UprightFont = * ExtraLight,
    BoldFont = * Bold,
    ItalicFont = AR PL UKai CN,       % 中文斜体通常用楷体替代
    BoldItalicFont = AR PL UKai CN    % 中文粗斜体
]
% 设置 CJK 无衬线字体（对应 \sffamily）为思源黑体 Light
\setCJKsansfont{Source Han Sans SC}[
    UprightFont = * ExtraLight,
    BoldFont = * Bold
]
% 设置 CJK 等宽字体（对应 \ttfamily）
\setCJKmonofont{FangSong}
% endregion
% endregion

% region tcolorbox 风格设置
% region 分页盒子的风格
\tcbset{
    pages/.style 2 args={
        blankest,
        breakable,
        break at=#1, % 11个公式一页高度（使用实际测量值）
        width=\textwidth/#2,  % 宽度固定为所有公式宽度
        before={\parindent0pt},
        size=tight,
        halign=center,
        height fixed for=all,
        reset and store to box array}}
% endregion
% region 定高盒子的风格
\tcbset{
    fixedheight/.style={
        blanker,height=#1,width=\linewidth,halign=center,valign=center}} % 11.3\gfllheight+2pt
% endregion
% region FLL 公式示意图盒子的风格
\tcbset{
    fll/.style={
        enhanced,
        colback=myblue,
        colframe=black!70,
        nobeforeafter,
        before skip=0pt,
        before={\parindent0pt},  % 去掉首行缩进
        after skip=0pt,
        borderline={1mm}{-1pt}{black!70},
        arc=5mm,
        boxsep=0pt,
        % 关键：使用tabularx属性直接创建n列表格（默认顶部对齐）
        tabularx={X*{#1}{|X}},  % n列等宽，顶部对齐，带分隔线
        width=\linewidth,}}  % 宽度占满整行
% endregion
% region 标题无边框设计
\tcbset{
    title without border/.style 2 args={
        empty,
        bottomtitle=1ex,
        grow to left by=-2em,
        frame code={\path[draw=#1,fill=#2,rounded corners=5mm,line width=1mm] (interior.south west) rectangle (interior.north east);},
        fonttitle=\sffamily\heiti\bfseries\Large,coltitle=#1
    }
}
% endregion
% endregion
\ExplSyntaxOn
% region 展开函数自定义变种
\exp_args_generate:n {eN}
\exp_args_generate:n {VN}
% endregion

% region 定义 MultiColsBox 盒子
\int_new:N \l_multicolsbox_columns_int
\int_new:N \l_multicolsbox_columns_minusone_int
\int_new:N \l_multicolsbox_rows_int
\box_new:N \l_multicolsbox_box_box
\seq_new:N \l_multicolsbox_boxarray_seq
\seq_new:N \l_multicolsbox_keys_seq
\seq_new:N \l_multicolsbox_vals_seq
\dim_new:N \l_multicolsbox_height_dim
\dim_new:N \l_multicolsbox_width_dim
\dim_new:N \l_multicolsbox_total_width_dim
\fp_new:N \l_multicolsbox_rows_fp
\tl_new:N \l_multicolsbox_coffin_tl
\coffin_new:N \l_multicolsbox_coffin
\bool_new:N \l_multicolsbox_ifcoffin_bool
\bool_set_false:N \l_multicolsbox_ifcoffin_bool

\keys_define:nn { MultiColsBox }
{
    columns .code:n = {
                        \int_set:Nn \l_multicolsbox_columns_int {#1}
                        \int_set:Nn \l_multicolsbox_columns_minusone_int {\int_eval:n {#1 - 1}}
                    },
    columns .default:n = 4,
    rows .code:n =  {
                        \int_set:Nn \l_multicolsbox_rows_int {#1}
                        \fp_set:Nn \l_multicolsbox_rows_fp {\fp_eval:n {1.02*#1}}
                    },
    rows .value_required:n = true,
    box .code:n =   {
                        \box_set_eq:NN \l_multicolsbox_box_box #1
                        \dim_set:Nn \l_multicolsbox_height_dim {\box_ht_plus_dp:N \l_multicolsbox_box_box}
                        \dim_set:Nn \l_multicolsbox_width_dim {\box_wd:N \l_multicolsbox_box_box}
                    },
    box .value_required:n = true,
    coffin .code:n = {
                        \tl_set:Nn \l_multicolsbox_coffin_tl {#1}
                        \bool_set_true:N \l_multicolsbox_ifcoffin_bool
                    },
    width .dim_set:N = \l_multicolsbox_total_width_dim,
    width .default:n = {\linewidth},
    % unknown .code:n = {
    %                     % 使用原始键名（包含空格）重新构建键值对
    %                     \seq_put_right:NV \l_multicolsbox_keys_seq \l_keys_key_tl
    %                     \seq_put_right:NV \l_multicolsbox_vals_seq #1
    %                 }
}
% \NewDocumentCommand{\addeq}{ m m }{#1=#2,}
\NewDocumentEnvironment{MultiColsBox}{ O{} O{} }
{
    \keys_set:nn {MultiColsBox} {#1}
    \begin{tcolorbox}[pages={\dim_eval:n {\fp_use:N \l_multicolsbox_rows_fp \l_multicolsbox_height_dim}}{\int_use:N \l_multicolsbox_columns_int}]
}
{
    \end{tcolorbox}
    \int_step_inline:nn {\int_use:N \l_multicolsbox_columns_int} {
        \seq_put_right:Nn \l_multicolsbox_boxarray_seq {
            \vbox{\vspace{1ex}\useboxarray{##1}}
        }
    }
    \vcoffin_set:Nnn \l_multicolsbox_coffin {\l_multicolsbox_total_width_dim} {
    \begin{tcolorbox}[fixedheight={\dim_eval:n {\int_use:N \l_multicolsbox_rows_int \l_multicolsbox_height_dim}}]
        \begin{tcolorbox}[fll={\int_use:N \l_multicolsbox_columns_minusone_int},#2]
            % 直接使用储存的页面，顶部对齐
            \seq_use:Nn \l_multicolsbox_boxarray_seq { & }
        \end{tcolorbox}
    \end{tcolorbox}}
    \bool_log:N \l_multicolsbox_ifcoffin_bool
    \tl_log:N \l_multicolsbox_coffin_tl
    \bool_if:NTF \l_multicolsbox_ifcoffin_bool
        {\exp_args:NeN \coffin_gset_eq:cN {\l_multicolsbox_coffin_tl} \l_multicolsbox_coffin}
        {\coffin_typeset:Nnnnn \l_multicolsbox_coffin {H} {l} {0pt} {0pt}}
    \bool_set_false:N \l_multicolsbox_ifcoffin_bool
    % \seq_log:N \l_multicolsbox_keyvals_seq
    \seq_clear:N \l_multicolsbox_keyvals_seq
}
% endregion

% region 一次性设置序列（用逗号分隔）
% 定义序列（变量名中不能有数字）
\seq_new:N \g_cube_notation_all_seq
\seq_new:N \g_cube_fll_all_seq
\seq_new:N \g_cube_oll_all_seq
\seq_new:N \g_cube_pll_all_seq

% region \g_cube_fll_all_seq 设置魔方公式的各种记号解释
\seq_set_from_clist:Nn \g_cube_notation_all_seq {R,R',R2,r,L,l,U,u,D,d,F,f,B,b,S,M,E,x,y,z}
% endregion

% region \g_cube_fll_all_seq 设置 F2L 阶段 41 个公式
\seq_set_from_clist:Nn \g_cube_fll_all_seq {
    {---rb}{rb---}{---------}{(R\blue{U'}R'U)y'\\(R'U2R\red{U'}\blue{U'})\\(R'UR)}, % 1
    {-b--b}{r----}{-------r-}{(UR\blue{U'}R'\blue{U'})\\y'(R'UR)}, % 2
    {----b}{r--r-}{-----b---}{(R'\red{F'}RU)\\(R\blue{U'}R'F)}, % 3
    {---bw}{br---}{---------}{(RUR'\blue{U'})\\(RU2R'\blue{U'})\\(RUR')}, % 4
    {---rw}{bb---}{---------}{(RF)U\\(R\blue{U'}R'\red{F'})\\(\blue{U'}R')或者\\(R\blue{U'}R'\blue{U'})\\(R\blue{U'}R'\blue{U})y'R'\blue{U'}R}, % 5
    {-b--w}{b----}{-------r-}{y'(R'\blue{U'}RU)\\(R'\blue{U'}R)}, % 6
    {----w}{b--r-}{-----b---}{(R\blue{U'}R'U)\\(R\blue{U'}R')}, % 7
    {---br}{wr---}{---------}{(R\blue{U'}R'U)\\(R\red{U'}\blue{U'}R'U)\\(R\blue{U'}R')}, % 8
    {---rr}{wb---}{---------}{(RU)F\\(RUR'\blue{U'})\\\red{F'}R'}, % 9
    {-b--r}{w----}{-------r-}{y'(R'UR\blue{U'})\\(R'UR)}, % 10
    {----r}{w--r-}{-----b---}{(RUR'\blue{U'})\\(RUR')}, % 11
    {--rb-}{-rb--}{--------w}{(RUR'\blue{U'})2\\(RUR')}, % 12
    {--rr-}{-bb--}{--------w}{(R\blue{U'}R'U)\\y'(R'UR)\\或者\\y'R'UR'FRF'R}, % 13
    {-br--}{--b--}{-------rw}{y'(R'U2)\\(RUR'\blue{U'}R)}, % 14
    {--r--}{--b--}{---r----w}{y'\blue{U'}\\(R'U2)\\(R\blue{U'}R'UR)}, % 15
    {--r--}{--b--}{-r------w}{y'(R'UR\red{U'}\blue{U'})\\(R'\blue{U'}R)\\或者\\\blue{F'}L'U2LF}, % 16
    {--r--}{--bb-}{-----r--w}{FU\\(R\blue{U'}R'\red{F'})\\(R\blue{U'}R')}, % 17
    {-rr--}{--b--}{-------bw}{U(R\blue{U'}R'\blue{U'})\\(R\blue{U'}R'U)\\(R\blue{U'}R')}, % 18
    {--r--}{--b--}{---b----w}{(R\blue{U'}R'U2)\\(RUR')\\或者\\yFRU2R'\red{F'}}, % 19
    {--r--}{--b--}{-b------w}{U(R\red{U'}\blue{U'})\\(R'UR\blue{U'}R')}, % 20
    {--r--}{--br-}{-----b--w}{(R\red{U'}\blue{U'})\\(R'\blue{U'}RUR')}, % 21
    {--wb-}{-rr--}{--------b}{\blue{U'}(R\blue{U'}R'U2)\\(R\blue{U'}R')}, % 22
    {--wr-}{-br--}{--------b}{\blue{U'}(RUR')\\y'(UR'\blue{U'}R)}, % 23
    {-bw--}{--r--}{-------rb}{y'U\\(R'UR\blue{U'})\\(R'\blue{U'}R)}, % 24
    {--w--}{--r--}{---r----b}{y'(R'\blue{U'}R)}, % 25
    {--w--}{--r--}{-r------b}{y'U\\(R'\blue{U'}R\blue{U'})\\(R'\blue{U'}R)}, % 26
    {--w--}{--rb-}{-----r--b}{y'(R\red{U'}\blue{U'})\\(R'2\blue{U'})\\(R2\blue{U'}R')}, % 27
    {-rw--}{--r--}{-------bb}{(lU)\\(r\blue{U'}r'\blue{U'})l'}, % 28
    {--w--}{--r--}{---b----b}{\blue{U'}(R\red{U'}\blue{U'})\\(R'U2)(R\blue{U'}R')\\或者\\y'rU2R'2\\U'R2U'r'}, % 29
    {--w--}{--r--}{-b------b}{\blue{U'}(RUR'\blue{U'})\\(R\red{U'}\blue{U'}R')}, % 30
    {--w--}{--rr-}{-----b--b}{UR\blue{U'}R'}, % 31
    {--bb-}{-rw--}{--------r}{\blue{U'}(R\red{U'}\blue{U'}R'U)\\(RUR')}, % 32
    {--br-}{-bw--}{--------r}{y'U(R'\blue{U'}R)\\\green{d'}(RUR')\\或者\\U(F'U'F)\\U'(RUR')}, % 33
    {-bb--}{--w--}{-------rr}{y'\blue{U'}\\(R'UR)}, % 34
    {--b--}{--w--}{---r----r}{y'U\\(R'\blue{U'}R\red{U'}\blue{U'})\\(R'UR)}, % 35
    {--b--}{--w--}{-r------r}{y'U\\(R'U2R\red{U'}\blue{U'})\\(R'UR)}, % 36
    {--b--}{--wb-}{-----r--r}{(R\blue{U'}R'U2)\\y'(R'\blue{U'}R)}, % 37
    {-rb--}{--w--}{-------br}{(R\blue{U'}R'U)\\(R\blue{U'}R'U2)\\(R\blue{U'}R')}, % 38
    {--b--}{--w--}{---b----r}{\blue{U'}(RUR'U)\\(RUR')}, % 39
    {--b--}{--w--}{-b------r}{(R U R')}, % 40
    {--b--}{--wr-}{-----b--r}{\blue{U'}(R\blue{U'}R'U)\\(RUR')} % 41
}
% endregion

% region \g_cube_oll_all_seq 设置 OLL 阶段 57 个公式
\seq_set_from_clist:Nn \g_cube_oll_all_seq {
    {LURL-RLDR}{(R\red{U'}\blue{U'})\\(R'2FR\red{F'})\\U2(R'FR\red{F'})}, % 01
    {LUUL-RLDD}{F(RUR'\blue{U'})\red{F'}\\f(RUR'\blue{U'})\red{f'}}, % 02
    {UURL-RLD-}{f(RUR'\blue{U'})\red{f'}\\\blue{U'}\\F(RUR'\blue{U'})\red{F'}}, % 03
    {LU-L-RDDR}{f(RUR'\blue{U'})\red{f'}\\\red{f'}U\\F(RUR'\blue{U'})\red{F'}}, % 04
    {UURL--L--}{(r'U2)\\(RUR'U)r\\或者\\y2l'U2L\\UL'Ul}, % 05
    {L--L--DDR}{(r\red{U'}U')\\(R'\blue{U'}R\blue{U'}r')\\或者\\y2lU2L'U'\\LU'l'}, % 06
    {U-R--R-DD}{(rUR'U)\\(R\red{U'}\blue{U'}r')\\或者\\y2lUL'U\\LU2l'}, % 07
    {-UU--RD-R}{(r'\blue{U'}R\blue{U'})\\(R'U2r)\\或者\\y2l'U'LU'\\L'U2l}, % 08
    {L-U--RDD-}{(RUR'\blue{U'})\\(R'F)\\(R2UR'\blue{U'})\red{F'}}, % 09
    {UU---RL-D}{(RUR'\green{U})\\(R'FR\red{F'})\\(R\red{U'}\blue{U'}R')}, % 10
    {UURL----D}{(r'R2UR'U)\\(R\red{U'}\blue{U'}R'U)\\\green{M'}}, % 11
    {--UL--DDR}{\green{M'}(R'\blue{U'}R\blue{U'})\\(R'U2R\blue{U'})M}, % 12
    {UU----LDD}{f(RU)\\(R2\blue{U'})\\(R'UR\blue{U'})\green{f'}}, % 13
    {LUU---DD-}{(R'F)\\(RUR'\red{F'}R)\\(F\blue{U'}\red{F'})}, % 14
    {UUR---LD-}{(r'\blue{U'}r)\\(R'\blue{U'}RU)\\(r'Ur)}, % 15
    {LU----DDR}{(rUr')\\(RUR'\blue{U'})\\(r\blue{U'}r')}, % 16
    {-UUL-RLD-}{(RUR'U)\\(R'FRF')\\U2(R'FRF')}, % 17
    {-U-L-RDDD}{(rUR'U)\\(R\red{U'}\blue{U'}r')\\(r'\blue{U'}R\blue{U'})\\(R'U2r)}, % 18
    {-U-L-RLDR}{(r'RU)\\(RUR'\blue{U'}r)\\(R'2FR\red{F'})}, % 19
    {-U-L-R-D-}{(rUR'\blue{U'})\\\green{M'2}U\\(R\blue{U'}R'\blue{U'})\green{M'}}, % 20
    {U-U---D-D}{(R\red{U'}\blue{U'})\\(R'\blue{U'}RUR'\blue{U'})\\(R\blue{U'}R')}, % 21
    {L-U---L-D}{(R\red{U'}\blue{U'})\\(R'2\blue{U'})(R2\blue{U'})\\(R'2\blue{U'})(\red{U'}R)}, % 22
    {U-U------}{(R'2\green{D'})\\(R\red{U'}\blue{U'}R'\green{D})\\(R\red{U'}\blue{U'}R)}, % 23
    {U-----D--}{(rUR'\blue{U'})\\(r'FR\red{F'})}, % 24
    {L-------D}{\red{F'}(rUR'\blue{U'})\\(r'FR)}, % 25
    {L-----D-R}{(R\red{U'}\blue{U'})\\(R'\blue{U'}R\blue{U'}R')\\或者\\y'R'U'RU'\\RU2R'}, % 26
    {U-R-----D}{(RUR'\green{U})\\(R\red{U'}\blue{U'}R')\\或者\\y'R'U2R\\UR'UR}, % 27
    {-----R-D-}{(rUR'\blue{U'})\\(r'RU)\\(R\blue{U'}R')}, % 28
    {U----RDD-}{(RUR'\blue{U'})\\(R\blue{U'}R'\red{F'}\blue{U'}F)\\(RUR')}, % 29
    {-U-L--L-R}{f(RU)\\(R2\blue{U'}R'U)\\(R2\blue{U'}R')\red{f'}}, % 30
    {-UU--R--D}{R'FRUR'\\U'F2UFR}, % 31
    {UU-L--D--}{(RU)(\red{B'}\blue{U'})\\(R'URBR')\\或者\\S(RUR'U')\\(R'FRf')}, % 32
    {UU----DD-}{(RUR'\blue{U'})\\(R'FR\red{F'})}, % 33
    {LUR----D-}{(RUR2\blue{U'})\\(R'F)\\(RUR\blue{U'}\green{F'})}, % 34
    {-URL--D--}{(R\red{U'}\blue{U'})\\(R'2FR\red{F'})\\(R\red{U'}\blue{U'}R')}, % 35
    {-UR--RD--}{(R'\blue{U'}R\blue{U'})\\(R'URU)\\(l\blue{U'}R'U)}, % 36
    {--R--RDD-}{F(R\blue{U'}R'\blue{U'})\\(RUR'\red{F'})}, % 37
    {U----R-DR}{(RUR'U)\\(R\blue{U'}R'\blue{U'})\\(R'FR\red{F'})}, % 38
    {LU-----DD}{(RUR'\red{F'}\blue{U'}F)\\U(RU2R')}, % 39
    {-UU---LD-}{(R'F)\\(RUR'\blue{U'})\\\red{F'}(UR)}, % 40
    {U-U--R-D-}{(RUR'\green{U})\\(R\red{U'}\blue{U'}R')F\\(RUR'\blue{U'})\red{F'}}, % 41
    {-U---RD-D}{(R'\blue{U'}R\blue{U'})\\(R'U2R)F\\(RUR'\blue{U'})\red{F'}}, % 42
    {-UR--R--R}{(\red{B'}\blue{U'})\\(R'URB)\\或者\\f'(L'U'LU)f}, % 43
    {LU-L--L--}{f(RUR'\blue{U'})\red{f'}}, % 44
    {LU----LD-}{F(RUR'\blue{U'})\red{F'}}, % 45
    {--RL-R--R}{(R'\blue{U'})\\(R'FR\red{F'})\\(UR)}, % 46
    {U-RL--DDR}{\red{b'}(\blue{U'}R'UR)2b\\或者\\F'(L'U'LU)2F}, % 47
    {L-U--RLDD}{F(RUR'\blue{U'})2\red{F'}}, % 48
    {UUR--RD-R}{(R\red{B'})\\(R'2FR2B)\\(R'2\red{F'}R)}, % 49
    {LUUL--L-D}{(r'U)\\(r2\blue{U'}r'2\blue{U'})\\(r2Ur')}, % 50
    {LUU---LDD}{f(RUR'\blue{U'})2\red{f'}}, % 51
    {L-UL-RL-D}{(R'\red{F'}\blue{U'}F\blue{U'})\\(RUR'\green{U}R)}, % 52
    {UUU--RD-D}{(r'U2)\\(RUR'\blue{U'})\\(RUR'Ur)}, % 53
    {U-U--RDDD}{(r\red{U'}\blue{U'})\\(R'\blue{U'}RU)\\(R'\blue{U'}R\blue{U'}r')}, % 54
    {UUU---DDD}{(r\red{U'}\blue{U'}R'\blue{U'})\\(r'R2UR'\blue{U'})\\(r\blue{U'}r')}, % 55
    {-U-----D-}{(RUR'\blue{U'})\\\green{M'}(UR\blue{U'}r')}, % 57
    {LUR---LDR}{(rUr')\\(UR\blue{U'}R')2\\(r\blue{U'}r')\\或者\\(fRUR'U'F')\\(RUR'U')\\(R'F'Rf')}, % 56
}
% endregion

% region \g_cube_pll_all_seq 设置 PLL 阶段 21 个公式
\seq_set_from_clist:Nn \g_cube_pll_all_seq {
    {l-l,uuu,r-r,d-d}{4-8-6-4}{Ua}{\green{M'2}U\green{M}U2\\\green{M'}U\green{M'2}\\或者\\(RU'R)(UR)2\\(U'R'U'R'2)}, % 01
    {l-l,uuu,r-r,d-d}{4-6-8-4}{Ub}{\green{M'2}\blue{U'}\green{M}U2\\\green{M'}\blue{U'}\green{M'2}\\或者\\(R2U)(RU)\\(R'U')2(R'UR')}, % 02
    {l-l,u-u,r-r,d-d}{2=8,4=6}{H} {\green{M'2}U\green{M'2}\\U2\\\green{M'2}U\green{M'2}}, % 03
    {l-l,u-u,r-r,d-d}{2=6,4=8}{Z} {\green{M'}U\\(\green{M'2}U)2\\\green{M'}U2\\\green{M'2}\blue{U'}}, % 04
    {-ll,uu,-r,-d   }{3-9-7-3}{Aa}{x'R2\green{D2}\\(R'\blue{U'}R)\green{D2}\\(R'UR')}, % 05
    {-ll,uu,-r,-d   }{3-7-9-3}{Ab}{x'(RU'R)D2\\(R'UR)D2R'2}, % 06
    {-l,-u,-r,-d    }{1=7,3=9}{E} {x'(R\blue{U'}R'\green{D})\\(RUR'\green{D'})\\(RUR'\green{D})\\(R\blue{U'}R'\green{D'})}, % 07
    {l-l,uu,-dd     }{3=9,4=6}{T} {(RUR'\blue{U'})\\(R'F)\\(R2\blue{U'}R'\blue{U'})\\(RUR'\red{F'})}, % 08
    {lll,u,-r,--d   }{2=8,3=9}{F} {(R'\blue{U'}\red{F'})\\(RUR'\blue{U'})\\(R'F)\\(R2\blue{U'}R'\blue{U'})\\(RUR'UR)}, % 09
    {ll,--u,r,-dd   }{1=9,2=6}{V} {R'URU'R'f'\\(U'RU'2R')\\(U'RU'R')fR}, % 10
    {rr,-dd         }{1=9,2=4}{Y} {F(R\blue{U'}R'\blue{U'})\\(RUR'\red{F'})\\(RUR'\blue{U'})\\(R'FR\red{F'})}, % 11
    {-ll,uuu,r      }{6=8,7=9}{Ja}{z(\blue{U'}R\red{D'})\\(R2UR'\blue{U'})\\(R2U)\\(\red{D}R')}, % 12
    {lll,uu,--d     }{3=9,6=8}{Jb}{(RUR'\blue{F'})\\(RUR'\blue{U'})\\(R'F)\\(R2\blue{U'}R'\blue{U'})}, % 13
    {ll,-u,--r,d-d  }{1=3,6=8}{Rb}{(R'U2)\\(R\red{U'}\blue{U'})\\(R'FRUR'\blue{U'})\\(R'\red{F'}R2\blue{U'})}, % 14
    {l-l,u--,-r,-dd }{2=4,3=9}{Ra}{(R\blue{U'}R'\blue{U'})\\(RUR\green{D})\\(R'\blue{U'}R\green{D'})\\(R'U2R'\blue{U'})}, % 15
    {l-l,-uu        }{1-7-9-1,4-6-8-4}{Gc}{(R'2\blue{u'}R\red{U'}R)\\(UR'u)\\(R2fR'\red{f'})}, % 16
    {l-l,-rr        }{1-3-7-1,2-4-8-2}{Gd}{(RUR')y'\\(R2\red{u'}R\blue{U'})\\(R'UR'uR2)}, % 17
    {l-l,dd         }{1-3-7-1,2-4-6-2}{Ga}{R2UR'UR'\\U'RU'R2\\(U'D)R'UR\\(D'U)}, % 18
    {l-l,rr         }{1-7-9-1,2-8-4-2}{Gb}{(R'd'F)\\(R2u)\\(R'U)\\(R\red{U'}R\blue{u'}R'2)}, % 19
    {l,-uu,r,-dd    }{1=9,4=6}{Nb}{(R'UR\blue{U'})\\(R'\red{F'}\blue{U'})\\(FRUR'F)\\(R'\red{F'}R\blue{U'}R)}, % 20
    {--l,uu,--r,dd  }{3=7,4=6}{Na}{(RUR'\green{U})\\(RUR'\blue{F'})\\(RUR'\blue{U'})(R'F)\\(R2\blue{U'}R'U2)\\(R\blue{U'}R')}, % 21
}
% endregion
% endregion
\ExplSyntaxOff % 启用LaTeX3语法

% \view{f} 正面斜 45度视角
\begin{document}

% region \PageOne
\NewCoffin \PageOne
\SetHorizontalCoffin \PageOne {}

% region \TitleCoffin 标题
\NewCoffin \TitleCoffin
\NewCoffin \LogoCoffin
\NewCoffin \TitleTextCoffin
\SetHorizontalCoffin \TitleCoffin {}
\SetHorizontalCoffin \LogoCoffin {\CubeLogo}
\SetVerticalCoffin \TitleTextCoffin {8cm} {{
    \raggedright\bfseries\sffamily\heiti\huge\color{title}
    三阶魔方{\Huge\bfseries}~CFOP~教程\\[1ex]
    $3\times3$ CFOP TUTORIAL
}}
\JoinCoffins \TitleCoffin [l,vc] \LogoCoffin [l,vc]
\JoinCoffins \TitleCoffin [\LogoCoffin-r,\LogoCoffin-vc] \TitleTextCoffin [l,vc] (2em,-1ex)
% endregion

% 把 \TitleCoffin 安装到 \PageOne 里
\JoinCoffins \PageOne [l,t] \TitleCoffin [l,t]

% region \StepsBoxCoffin CFOP 复原步骤
\newlength{\stepsgap}
\setlength{\stepsgap}{.5em}
\newlength{\stepsunit}
\setlength{\stepsunit}{1.5cm}
\NewCoffin \StepsCoffin
\SetHorizontalCoffin \StepsCoffin {}
\NewCoffin \ShortArrow
\SetHorizontalCoffin \ShortArrow {\tikz{\draw[orange,-{Stealth[length=1.5cm,width=1.5cm,inset=0pt]},line width=6mm] (0,0)--(2.5,0);}}
\NewCoffin \CrossStepCoffin
\SetHorizontalCoffin \CrossStepCoffin {\view{r}\tikz{\Cube{-r--r----}{-b--b----}{-w-www-w-}[\stepsunit]}}
\NewCoffin \FLLStepCoffin
\SetHorizontalCoffin \FLLStepCoffin {\tikz{\Cube{---bb}{rr---}{---------}[\stepsunit]}}
\NewCoffin \OLLStepCoffin
\SetHorizontalCoffin \OLLStepCoffin {\tikz{\Cube{---bb}{rr---}{yyyyyyyyy}[\stepsunit]}}
\NewCoffin \PLLStepCoffin
\SetHorizontalCoffin \PLLStepCoffin {\tikz{\Cube{bbbbb}{rrrrr}{yyyyyyyyy}[\stepsunit]}}

\JoinCoffins \StepsCoffin[hc,vc] \CrossStepCoffin[l,vc]
\JoinCoffins \StepsCoffin[\CrossStepCoffin-r,\CrossStepCoffin-vc] \ShortArrow[l,vc] (\stepsgap,0pt)
\JoinCoffins \StepsCoffin[\ShortArrow-r,\ShortArrow-vc] \FLLStepCoffin[l,vc] (\stepsgap,0pt)
\JoinCoffins \StepsCoffin[\FLLStepCoffin-r,\FLLStepCoffin-vc] \ShortArrow[l,vc] (\stepsgap,0pt)
\JoinCoffins \StepsCoffin[\ShortArrow-r,\ShortArrow-vc] \OLLStepCoffin[l,vc] (\stepsgap,0pt)
\JoinCoffins \StepsCoffin[\OLLStepCoffin-r,\OLLStepCoffin-vc] \ShortArrow[l,vc] (\stepsgap,0pt)
\JoinCoffins \StepsCoffin[\ShortArrow-r,\ShortArrow-vc] \PLLStepCoffin[l,vc] (\stepsgap,0pt)

\NewCoffin \CrossText
\NewCoffin \FLLText
\NewCoffin \OLLText
\NewCoffin \PLLText
\SetVerticalCoffin \CrossText {2.2\stepsunit} {
    \centering\sffamily\heiti\Large
    \blue{\huge\bfseries CROSS}\\底层十字}
\SetVerticalCoffin \FLLText {2.2\stepsunit} {
    \centering\sffamily\heiti\Large
    \blue{\huge\bfseries F2L}\\前两层归位}
\SetVerticalCoffin \OLLText {2.2\stepsunit} {
    \centering\sffamily\heiti\Large
    \blue{\huge\bfseries OLL}\\最后一层翻色}
\SetVerticalCoffin \PLLText {2.2\stepsunit} {
    \centering\sffamily\heiti\Large
    \blue{\huge\bfseries PLL}\\最后一层归位}

\ExplSyntaxOn
\clist_map_variable:nNn {\CrossText,\FLLText,\OLLText,\PLLText} \l_tmpa_tl
    {
        \tl_log:N \l_tmpa_tl
        \expandafter\ScaleCoffin\l_tmpa_tl {3} {3}}
\ExplSyntaxOff

\JoinCoffins \StepsCoffin[\CrossStepCoffin-b,\CrossStepCoffin-hc] \CrossText[t,hc] (0pt, -5\stepsgap)
\JoinCoffins \StepsCoffin[\FLLStepCoffin-b,\FLLStepCoffin-hc] \FLLText[t,hc] (0pt, -5\stepsgap)
\JoinCoffins \StepsCoffin[\OLLStepCoffin-b,\OLLStepCoffin-hc] \OLLText[t,hc] (0pt, -5\stepsgap)
\JoinCoffins \StepsCoffin[\PLLStepCoffin-b,\PLLStepCoffin-hc] \PLLText[t,hc] (0pt, -5\stepsgap)

\ScaleCoffin \StepsCoffin {0.25} {0.25}
\NewCoffin \StepsBoxCoffin
\SetHorizontalCoffin \StepsBoxCoffin{
    \tcbox[title without border={orange}{steps},title={CFOP 复原步骤}]{\TypesetCoffin \StepsCoffin}}
% endregion

% 把 \StepsBoxCoffin 安装到 \PageOne 里
\JoinCoffins \PageOne [\LogoCoffin-l,\LogoCoffin-b] \StepsBoxCoffin [l,t] (.5em,-1ex)

% region CubeNotation 盒子
\ExplSyntaxOn
\NewCoffin \NotationCoffin
\SetHorizontalCoffin \NotationCoffin {
\begin{tcolorbox}[
    title~without~border={nicered}{notation},
    title=魔方转动公式符号及示意图,
    % width=\textwidth-\CoffinWidth\PageOne-1cm,
    % width=7\gnotationwidth,
    width=\CoffinWidth\PageOne-1.1cm,
    height=\textheight-\CoffinTotalHeight\PageOne,
    ]
    \begin{tcbraster}[
        enhanced,
        size=small,
        attach~boxed~title~to~top~left,
        raster~columns=5,
        raster~equal~height,
        raster~row~skip=-1ex,
        % raster~row~skip=(\textheight-4\gnotationheight-3cm)/3,
        ]
        \seq_map_variable:NNn \g_cube_notation_all_seq \l_tmpa_tl
            {
                \begin{tcolorbox}[blankest,size=tight]
                    \exp_args:No\CubeNotation{\l_tmpa_tl}
                \end{tcolorbox}}
    \end{tcbraster}
\end{tcolorbox}
}
\ExplSyntaxOff
% endregion

% 把 \NotationCoffin 安装到 \PageOne 里
\JoinCoffins \PageOne [\StepsBoxCoffin-b,\StepsBoxCoffin-l] \NotationCoffin [l,t] (.5em, -1ex)

% region F2L 公式部分
\NewCoffin \FLLCoffin
\begin{MultiColsBox}[
    columns=5,
    rows=10,
    box={\gcubefllbox},
    coffin={FLLCoffin},
    width=\textwidth-\CoffinWidth\PageOne+1cm,
    ]
\ExplSyntaxOn
\int_step_inline:nn {9} {\hbox{\rlap{\rule{0pt}{.98\gfllheight}}\rule{\gfllwidth}{0pt}}}
\seq_map_variable:NNn \g_cube_fll_all_seq \l_tmpa_tl {
    \hbox{\exp_after:wN\FLL\l_tmpa_tl}
}
\ExplSyntaxOff
\end{MultiColsBox}
% endregion

% 把 \FLLCoffin 安装到 \PageOne 里
\JoinCoffins \PageOne [r,b] \FLLCoffin [l,b] (4.16pt, .25cm)

% region F2L 及第一页的说明文字
\NewCoffin \FLLTextCoffin
\ExplSyntaxOn
\dim_log:N \gfllwidth
\dim_new:N \l_cube_fll_width_dim
\dim_set:Nn \l_cube_fll_width_dim {\gfllwidth - .5cm}
\let\lfllwidth=\l_cube_fll_width_dim
\ExplSyntaxOff
\SetVerticalCoffin \FLLTextCoffin {\lfllwidth} {
    \raggedright\sffamily\heiti
    \Huge\bfseries\color{orange}
    F2L\\
    \Large
    前两层归位\\[2cm]
    \addtolength{\textwidth}{-2cm}
    \footnotesize\mdseries
    \blue{蓝色字母：用左食指拨动}\\
    \red{红色字母：用右拇指拨动}\\
    \green{绿色字母：用其它指拨动}\\[4cm]
    \color{black}
    小写字母：双层转\\
    字母：顺时针转90°\\
    字母+'：逆时针转90°\\
    字母+2：重复两次\\[2cm]
    \hspace{-3pt}\raisebox{-2ex}{\tikz[radius=1cm,scale=0.25]{\fill[orange] (0,0) -- (0,2) -- (5,2) arc[start angle=90,end angle=-90] -- cycle node[white,anchor=west,above right=1.2pt]{温馨提示};}}\color{orange}\dotfill\\[2ex]
    *CROSS复原方法较多，本教程主要介绍后面3步的复原公式
}
% endregion

% 把 \FLLTextCoffin 安装到 \PageOne 里
\JoinCoffins \PageOne [\FLLCoffin-l,\FLLCoffin-t] \FLLTextCoffin [l,t] (5pt, -10pt)

% 输出 \PageOne 
\TypesetCoffin \PageOne (-2cm,.5cm)
% endregion

% region \PageTwo
\NewCoffin \PageTwo
\SetHorizontalCoffin \PageTwo {}

% region OLL 公式页面
\NewCoffin \OLLCoffin
\begin{MultiColsBox}[
    columns=9,
    rows=7,
    box={\gcubeollbox},
    coffin={OLLCoffin},
    width=\textwidth,
    ]
\ExplSyntaxOn
\seq_map_variable:NNn \g_cube_oll_all_seq \l_tmpa_tl {
    \hbox{\exp_after:wN\OLL\l_tmpa_tl}
}
\ExplSyntaxOff
\end{MultiColsBox}
% endregion

% 把 \OLLCoffin 安装到 \PageTwo 里
\JoinCoffins \PageTwo [l,t] \OLLCoffin [l,t] % (0pt, -1cm)
% \TypesetCoffin  \OLLCoffin

% region PLL 公式页面
\NewCoffin \PLLCoffin
\begin{MultiColsBox}[
    columns=7,
    rows=3,
    box={\gcubepllbox},
    coffin={PLLCoffin},
    width=\textwidth,][
    sharp corners=northeast,
    ]
\ExplSyntaxOn
\seq_map_variable:NNn \g_cube_pll_all_seq \l_tmpa_tl {
    \hbox{\exp_after:wN\PLL\l_tmpa_tl}
}
\ExplSyntaxOff
\end{MultiColsBox}
% endregion

% 把 \PLLCoffin 安装到 \PageTwo 里
\JoinCoffins \PageTwo [\OLLCoffin-l,\OLLCoffin-b] \PLLCoffin [l,t] (0pt, -2ex)

% region \OLLTextCoffin
\NewCoffin \OLLTextCoffin
\SetVerticalCoffin \OLLTextCoffin {3cm} {
    \raggedright\sffamily\heiti
    \Large\bfseries\color{nicered}
    OLL\\
    \large
    最后一层翻色\\[2cm]
    \footnotesize\mdseries\color{black}
    OLL和PLL是俯视图
    \renewcommand\descriptionlabel[1]{\textsf{#1：}\hspace{\labelsep}}
    \begin{description}[left=0pt,itemindent=-18pt,labelsep=0pt,rightmargin=1em]
        \color{nicered}
        \item[OLL] 蓝块和小蓝条表示黄色块\\[2cm]
        \color{niceblue}
        \item[PLL] 小黑条表示已归位的色块
    \end{description}
}
% endregion

% 把 \OLLTextCoffin 安装到 \PageTwo 里
\JoinCoffins \PageTwo [\OLLCoffin-r,\OLLCoffin-t] \OLLTextCoffin [r,t] (0cm, -4cm)

% region \CutoffCoffin 贴片框
\NewCoffin \CutoffCoffin
\SetHorizontalCoffin \CutoffCoffin{
    \tcbox[
    enhanced,
    width=90pt,
    arc=5mm,
    boxsep=0em,
    leftupper=2mm,
    rightupper=1mm,
    frame code={\draw[black!70,rounded corners=5mm,line width=1mm] (frame.south west) -- ([yshift=-5mm]frame.north west) -- ([yshift=-5mm]frame.north east) -- (frame.north east);},
    interior code={\fill[white] ([xshift=.5mm,yshift=-1.5mm]frame.south west) [rounded corners=4.5mm] -- ([xshift=.5mm,yshift=-5.5mm]frame.north west) [rounded corners=5.5mm] -- ([xshift=.5mm,yshift=-5.5mm]frame.north east) [sharp corners] -- ([xshift=.5mm]frame.north east) -- ([xshift=2mm]frame.north east) -- ([xshift=2mm,yshift=-1.5mm]frame.south east) -- cycle;}
    ]
    {\hbox{\rlap{\rule{0pt}{15mm}}\rule{29.5mm}{0pt}}}
}
% endregion

% 把 \CutoffCoffin 安装到 \PageTwo 里
\JoinCoffins \PageTwo [\OLLCoffin-r,\OLLCoffin-b] \CutoffCoffin [r,b] (3pt, -2.5pt)

% region \PLLTextCoffin
\NewCoffin \PLLTextCoffin
\SetHorizontalCoffin \PLLTextCoffin {
    \begin{tcolorbox}[
    enhanced,
    width=90pt,
    arc=5mm,
    boxsep=0em,
    leftupper=2mm,
    rightupper=1mm,
    frame code={\draw[black!70,fill=myblue,line width=1mm] (frame.south west) [rounded corners=5mm] -- (frame.north west) [rounded corners=5mm] -- (frame.north east) [sharp corners] -- (frame.south east);},
    interior code={\fill[myblue] ([yshift=-1mm]interior.south west |- frame.south west) [rounded corners=4.5mm] -- (interior.north west) [rounded corners=4.5mm] -- (interior.north east) [sharp corners] -- ([yshift=-1mm]interior.south east |- frame.south east) -- cycle;},
    ]
        \raggedright\sffamily\heiti
        \Large\bfseries\color{niceblue}
        PLL\\
        \large
        最后一层归位
    \end{tcolorbox}}
% endregion

% 把 \PLLTextCoffin 安装到 \PageTwo 里
\JoinCoffins \PageTwo [\PLLCoffin-r,\PLLCoffin-t] \PLLTextCoffin [r,b] (-.5pt, -2pt)

% % 旋转 180° 适合手机版华为打印
% \RotateCoffin \PageTwo {180}
% \TypesetCoffin \PageTwo [r,t] (-.85cm,1ex)

% 输出 \PageTwo 
\TypesetCoffin \PageTwo  (-.85cm,2cm)
% endregion

\end{document}